% Author: Alfredo Sánchez Alberca (asalber@ceu.es}
% !TEX program = xelatex
% !LTeX: language=es
\documentclass[aspectratio=169,10pt,t]{beamer}
%-------------------------------------------------------------------------------
% GENERAL PACKAGES
%-------------------------------------------------------------------------------
% Language
\usepackage{polyglossia}
\setmainlanguage{spanish}
% Maths
\usepackage{amsmath} % Math symbols and environments
\usepackage{amsfonts}
\usepackage{amssymb}
% Tables
\usepackage{array}
\usepackage{multirow}
\usepackage{booktabs}

% Graphics
\usepackage{graphicx}
\usepackage{tikz}
\usetikzlibrary{positioning}

% Colors
\definecolor{blueceu}{RGB}{0,164,227}
\definecolor{greenceu}{RGB}{194,205,24}
\definecolor{redceu}{RGB}{238,50,36}
\definecolor{purpleceu}{RGB}{169,78,145}
\definecolor{greyceu}{RGB}{117,117,97}
\definecolor{darkgrey}{RGB}{40,40,50}
\definecolor{softblueceu}{RGB}{193,225,246}
\setbeamercolor{structure}{fg=blueceu}
\setbeamercolor{normal text}{fg=darkgrey}
\hypersetup{colorlinks, urlcolor=purpleceu}

% Boxes
\usepackage[most]{tcolorbox}
\usepackage{setspace}
\newtcolorbox{datos}{
  enhanced,
  colback=blueceu!10, 
  colframe=blueceu, 
  fonttitle=\bfseries, 
  left=3pt, 
  right=3pt, 
  boxrule=0.5pt,
  code={\setstretch{1.2}},
  title={Datos},
}



%-------------------------------------------------------------------------------
% FONTS
%-------------------------------------------------------------------------------
\usepackage{fontspec}
\setmainfont[Ligatures=TeX]{TeX Gyre Pagella}
\usepackage{unicode-math}
\setmathfont[math-style=ISO, bold-style=ISO]{TeX Gyre Pagella Math}
% Creative common icons
\usepackage[
    type={CC},
    modifier={by-nc-sa},
    version={3.0},
    imagemodifier={-eu}
]{doclicense}

%-------------------------------------------------------------------------------
% CONFIGURATION
%-------------------------------------------------------------------------------
\setbeamersize{text margin left=.5cm, text margin right=.5cm} % Defines margin sizes
\beamertemplatenavigationsymbolsempty % Hide navitation bar
\usefonttheme[onlymath]{serif} % Math text in serif
\setbeamertemplate{blocks}[rounded] % Blocks with rounded corners
%\setbeamercolor{block title}{bg=RoyalBlue!10} % Color of block title
%\setbeamercolor{block body}{bg=RoyalBlue!10} % Color of block body

%-------------------------------------------------------------------------------
% COMMANDS
%-------------------------------------------------------------------------------
% \newcommand{\sen}{\operatorname{sen}}
% \newcommand{\tg}{\operatorname{tg}}
% \newcommand{\arcsen}{\operatorname{arcsen}}
% \newcommand{\arctg}{\operatorname{arctg}}

%-------------------------------------------------------------------------------
% DOCUMENT
%-------------------------------------------------------------------------------
\begin{document}
%---------------------------------------------------------------------SLIDE----
\begin{frame}[c]
\vspace{1.5cm}

\begin{center}
\structure{\LARGE {\textbf{Ejercicios de Cálculo}}}
\bigskip

\large
\begin{tabular}{rl}
Temas: & \structure{Derivadas en n variables}\\
Titulaciones: & \structure{Todas}
\end{tabular}

\bigskip
Alfredo Sánchez Alberca\\
\url{asalber@ceu.es}\\
\url{https://aprendeconalf.es}\\

\includegraphics[scale=0.2]{../img/logo_uspceu}

\bigskip
\doclicenseIcon
\end{center}
\end{frame}

%---------------------------------------------------------------------SLIDE----
\begin{frame}[c]
\Large
La función $T(x,y)=\ln(3xy+2x^2-y)$ da la temperatura de la superficie de una montaña en la latitud $x$ y longitud $y$. Unos montañeros están perdidos en la posición $(1,2)$ y corren el riesgo de morir congelados.

\begin{enumerate}
\item ¿En qué dirección deben moverse para evitar el riesgo de congelación lo más rápidamente posible?
\item Si se mueven en una dirección equivocada de manera que la longitud decrece la mitad de lo que aumenta la latitud, ¿aumentará o disminuirá el riesgo de hipotermia?
\item ¿En qué dirección deben moverse para que la temperatura permanezca constante?
\end{enumerate}
\end{frame}


%------------------------------------------------------------------SLIDE----
\begin{frame}
\begin{columns}
\begin{column}[T]{0.7\textwidth}
La función $T(x,y)=\ln(3xy+2x^2-y)$ da la temperatura de la superficie de una montaña en la latitud $x$ y longitud $y$. Unos montañeros están perdidos en la posición $(1,2)$ y corren el riesgo de morir congelados.

\begin{enumerate}
\item ¿En qué dirección deben moverse para evitar el riesgo de congelación lo más rápidamente posible?
\end{enumerate}
\end{column}
\quad
\begin{column}[T]{0.3\textwidth}
\begin{datos}
$T(x,y)=\ln(3xy+2x^2-y)$ \\
Punto $(1,2)$\\
\end{datos}
\end{column}
\end{columns}
\end{frame}


%------------------------------------------------------------------SLIDE----
\begin{frame}
\begin{columns}
\begin{column}[T]{0.75\textwidth}
\begin{enumerate}
\item[2.] Si se mueven en una dirección equivocada de manera que la longitud decrece la mitad de lo que aumenta la latitud, ¿aumentará o disminuirá el riesgo de hipotermia?
\end{enumerate}
\end{column}
\begin{column}[T]{0.25\textwidth}
\begin{datos}
$\nabla T(1,2) = \frac{1}{3}(5,1)$
\end{datos}
\end{column}
\end{columns}
\end{frame}

%------------------------------------------------------------------SLIDE----
\begin{frame}
\begin{columns}
\begin{column}[T]{0.75\textwidth}
\begin{enumerate}
\item[3.] ¿En qué dirección deben moverse para que la temperatura permanezca constante?
\end{enumerate}
\end{column}
\begin{column}[T]{0.25\textwidth}
\begin{datos}
$\nabla T(1,2) = \frac{1}{3}(5,1)$
\end{datos}
\end{column}
\end{columns}
\end{frame}

\end{document}